\chapter{Pendahuluan}

\section{Latar Belakang}

\noindent Koperasi merupakan pilar ekonomi kerakyatan Indonesia yang tertuang dalam Undang-Undang Dasar 1945 Pasal 33 ayat (1) sebagai bentuk usaha bersama berdasarkan asas kekeluargaan. Berdasarkan data Kementerian Koperasi dan UKM, Indonesia memiliki lebih dari 127.000 unit koperasi dengan aset gabungan melebihi 10 miliar dolar AS, menjadikan koperasi sebagai salah satu motor penggerak ekonomi nasional \cite{sailana2023pengaruh}. Namun, di balik angka yang besar tersebut, koperasi Indonesia menghadapi sejumlah permasalahan struktural yang menggerogoti kepercayaan anggota dan menghambat potensi pertumbuhannya.

\medskip
\noindent Permasalahan pertama adalah tata kelola yang opak. Pencatatan keuangan secara manual---melalui buku kas fisik dan \textit{spreadsheet}---membuka celah manipulasi data simpanan, pinjaman, dan Sisa Hasil Usaha (SHU) tanpa jejak audit yang memadai \cite{kartika2024tanggung}. Korupsi dana koperasi yang terjadi secara berulang---dari penggelapan simpanan hingga mark-up pinjaman fiktif---menunjukkan bahwa mekanisme pengawasan konvensional tidak cukup efektif. Permasalahan kedua adalah partisipasi anggota yang rendah. Rapat Anggota Tahunan (RAT) yang seharusnya menjadi forum pengambilan keputusan tertinggi seringkali hanya dihadiri oleh segelintir anggota aktif, sementara mayoritas anggota pasif karena tidak memiliki akses terhadap informasi terkini tentang kondisi koperasi \cite{novatiani2023pengaruh}. Permasalahan ketiga adalah distribusi SHU yang tidak transparan. Ketentuan Pasal 45 UU 25/1992 mengamanatkan pembagian SHU secara adil dan proporsional, namun dalam praktiknya, perhitungan dilakukan secara manual dan rentan terhadap kesalahan maupun manipulasi \cite{sailana2023pengaruh}.

\medskip
\noindent Teknologi blockchain dan \textit{Decentralized Autonomous Organization} (DAO) menawarkan solusi potensial terhadap permasalahan di atas. Blockchain menyediakan \textit{distributed ledger} yang imutabel---setiap transaksi tercatat secara permanen dan dapat diverifikasi oleh siapa pun tanpa bergantung pada otoritas pusat \cite{nakamoto2008bitcoin}. Ethereum memperluas konsep blockchain dengan memperkenalkan \textit{smart contract}---program yang berjalan di atas blockchain dan mengeksekusi aturan secara otomatis tanpa intervensi manusia \cite{buterin2014ethereum}. \textit{Smart contract} memungkinkan aturan bisnis dikodekan secara eksplisit dan dieksekusi secara deterministik, menciptakan transparansi dan akuntabilitas yang tidak mungkin dicapai oleh sistem manual.

\medskip
\noindent Namun, model DAO yang ada saat ini---seperti MakerDAO dan Compound---menggunakan mekanisme \textit{token-weighted voting}, di mana hak suara proporsional terhadap jumlah token yang dimiliki \cite{sun2022multiparty}. Model ini menciptakan dinamika plutokratis: anggota dengan kepemilikan token besar mendominasi keputusan tata kelola, sementara anggota kecil kehilangan suara secara efektif. Mekanisme ini bertentangan secara fundamental dengan prinsip koperasi yang diamanatkan oleh UU 25/1992 Pasal 22 ayat (1): setiap anggota memiliki satu suara tanpa memandang besarnya kontribusi modal. Kesenjangan antara model DAO \textit{token-weighted} dengan prinsip demokrasi koperasi menciptakan celah riset yang signifikan---diperlukan pendekatan baru yang menggabungkan transparansi dan otomatisasi blockchain dengan prinsip satu-anggota-satu-suara yang menjadi fondasi koperasi Indonesia \cite{sharma2024future}.

\medskip
\noindent Regulasi perkoperasian Indonesia---khususnya UU 25/1992 tentang Perkoperasian---mengandung 26 ketentuan yang bersifat operasional dan dapat diterjemahkan ke dalam aturan yang dapat dieksekusi secara digital. Ketentuan tersebut mencakup: keanggotaan sukarela dan terbuka (Pasal 5 ayat 1), pengelolaan demokratis (Pasal 5 ayat 2), satu anggota satu suara (Pasal 22), simpanan pokok, wajib, dan sukarela (Pasal 41), Rapat Anggota Tahunan (Pasal 26--27), pengawas dengan hak veto (Pasal 38--39), pembagian SHU (Pasal 45), dan laporan keuangan (Pasal 46--47) \cite{maryam2025pendirian}. Hingga saat ini, belum ada sistem yang secara sistematis memetakan seluruh ketentuan operasional UU 25/1992 ke dalam \textit{smart contract} yang dapat dieksekusi dan diverifikasi. Penelitian-penelitian sebelumnya tentang blockchain untuk koperasi---seperti Arisudhana dkk. (2025) yang membahas prinsip koperasi dalam konteks blockchain dan Sailana dkk. (2023) yang menganalisis simpanan dan SHU---masih bersifat konseptual dan parsial, belum menghasilkan implementasi on-chain yang komprehensif \cite{arisudhana2025pelatihan}.

\medskip
\noindent DChain (sebelumnya DiktiChain) merupakan jaringan blockchain konsorsium pendidikan tinggi Indonesia yang kompatibel dengan Ethereum Virtual Machine (EVM). DChain beroperasi pada Chain ID 17845 dan menyediakan infrastruktur blockchain yang dikelola oleh konsorsium perguruan tinggi, menawarkan lingkungan yang lebih terkendali dan sesuai dengan konteks Indonesia dibandingkan dengan \textit{public blockchain} seperti Ethereum mainnet \cite{wood2014ethereum}. Keberadaan DChain memberikan peluang bagi pengembangan sistem berbasis blockchain yang sesuai dengan kebutuhan institusi pendidikan dan pemerintahan di Indonesia, termasuk implementasi koperasi digital yang mematuhi regulasi nasional.

\medskip
\noindent Penelitian ini mengimplementasikan Lakomi---sebuah sistem koperasi berbasis blockchain yang memetakan 26 ketentuan operasional UU 25/1992 ke dalam empat \textit{smart contract}: LakomiToken (keanggotaan), LakomiVault (perbendaharaan dan SHU), LakomiGovern (tata kelola), dan LakomiLoans (pinjaman). Sistem di-deploy pada jaringan DChain dan diuji melalui 26 skenario pengujian fungsional yang memverifikasi kepatuhan terhadap setiap ketentuan. Pendekatan \textit{Design Science Research} (DSR) digunakan untuk memastikan bahwa \textit{artifact} yang dihasilkan tidak hanya berfungsi secara teknis tetapi juga memiliki justifikasi ilmiah yang kuat \cite{hevner2004design}. Lakomi menghadirkan arsitektur \textit{dual-track}: kontribusi finansial anggota memengaruhi kapasitas pinjaman (\textit{loan-to-value}), namun hak suara dalam tata kelola tetap setara---satu anggota, satu suara---sesuai amanat Pasal 22(1) UU 25/1992.

\section{Rumusan Masalah}

Berdasarkan latar belakang yang telah diuraikan, rumusan masalah dalam penelitian ini adalah:

\begin{enumerate}
	\item Bagaimana merancang dan mengimplementasikan \textit{smart contract} yang memetakan ketentuan UU 25/1992 ke dalam sistem koperasi berbasis blockchain pada jaringan DChain?
	\item Bagaimana memverifikasi kepatuhan 26 ketentuan UU 25/1992 yang telah dipetakan ke dalam fungsi \textit{smart contract} melalui pengujian fungsional?
	\item Bagaimana mengevaluasi kinerja dan kelayakan sistem melalui analisis biaya gas (\textit{gas analysis}) dan perbandingan dengan sistem koperasi tradisional serta DAO berbasis \textit{token-weighted voting}?
\end{enumerate}

\noindent\textbf{Justifikasi Urgensi.} Apabila sistem koperasi tidak mengadopsi teknologi blockchain, maka: (1) pencatatan keuangan tetap dilakukan secara manual dan rawan terhadap manipulasi tanpa jejak audit yang kredibel; (2) partisipasi anggota dalam tata kelola tetap rendah karena tidak ada mekanisme pemungutan suara yang transparan dan dapat diakses dari mana saja; (3) distribusi SHU tetap bergantung pada perhitungan manual yang tidak transparan; (4) pengawasan oleh pengawas koperasi tetap lemah karena tidak memiliki akses \textit{real-time} terhadap data keuangan; (5) kepatuhan terhadap UU 25/1992 tetap bergantung pada itikad baik pengurus tanpa verifikasi independen yang dapat diaudit publik; dan (6) Indonesia kehilangan peluang untuk menjadi pelopor dalam digitalisasi koperasi berbasis blockchain di tingkat nasional maupun internasional.

\section{Tujuan Penelitian}

Tujuan penelitian ini adalah:

\begin{enumerate}
	\item Merancang dan mengimplementasikan empat \textit{smart contract}---LakomiToken, LakomiVault, LakomiGovern, dan LakomiLoans---yang memetakan 26 ketentuan UU 25/1992 ke dalam sistem koperasi berbasis blockchain pada jaringan DChain.
	\item Memverifikasi kepatuhan 26 ketentuan UU 25/1992 melalui pengujian fungsional berbasis skenario yang mencakup seluruh siklus operasional koperasi: pendaftaran anggota, simpanan, pinjaman, tata kelola, pemilihan pengurus, pengawasan, distribusi SHU, dan RAT tahunan.
	\item Mengevaluasi kinerja dan kelayakan sistem melalui analisis biaya gas untuk tujuh operasi utama koperasi serta perbandingan komprehensif dengan sistem koperasi tradisional dan DAO berbasis \textit{token-weighted voting}.
\end{enumerate}

\section{Batasan Penelitian}

Penelitian ini memiliki batasan sebagai berikut:

\begin{enumerate}
	\item \textbf{Lingkungan Deployment.} Sistem di-deploy pada jaringan DChain \textit{testnet} (Chain ID 17845), sehingga evaluasi dilakukan dalam lingkungan pengembangan terkendali, bukan pada lingkungan produksi \textit{mainnet} dengan pengguna nyata.
	
	\item \textbf{Cakupan Regulasi.} Implementasi mencakup 26 ketentuan UU 25/1992 yang bersifat operasional dan dapat diterjemahkan ke dalam aturan yang dapat dieksekusi oleh \textit{smart contract}. Ketentuan yang bersifat administratif murni---seperti prosedur pendirian badan hukum koperasi, pembuatan Anggaran Dasar/Anggaran Rumah Tangga (AD/ART) oleh notaris, dan perizinan usaha---tidak termasuk dalam lingkup penelitian karena tidak dapat diotomatisasi melalui \textit{smart contract}.
	
	\item \textbf{Pengujian Keamanan.} Pengujian yang dilakukan terbatas pada pengujian fungsional untuk verifikasi kepatuhan terhadap UU 25/1992 dan analisis biaya gas. Audit keamanan formal---termasuk \textit{penetration testing}, verifikasi formal (\textit{formal verification}), dan analisis kerentanan menggunakan \textit{tool} seperti Slither atau Mythril---berada di luar lingkup penelitian ini.
	
	\item \textbf{Frontend.} Antarmuka pengguna (\textit{frontend}) dikembangkan menggunakan React 18 dan wagmi v2 sebagai prototipe untuk mendemonstrasikan interaksi dengan \textit{smart contract}. Optimasi \textit{production-grade}---termasuk \textit{server-side rendering}, \textit{caching}, dan \textit{responsive design} penuh---tidak dilakukan.
	
	\item \textbf{Manajemen Identitas.} Lima peran dalam sistem (PENGAWAS, APPROVER, MEMBERSHIP, TREASURER, GOVERN) diimplementasikan menggunakan akun Ethereum yang terpisah pada lingkungan pengujian. Integrasi dengan sistem identitas eksternal---seperti OAuth, \textit{Know Your Customer} (KYC), atau \textit{Single Sign-On}---berada di luar lingkup.
	
	\item \textbf{Interoperabilitas.} Sistem dirancang untuk beroperasi pada satu jaringan blockchain (DChain). Interoperabilitas \textit{cross-chain}---seperti transfer aset antar \textit{chain} atau integrasi dengan \textit{oracle} eksternal---tidak diimplementasikan.
	
	\item \textbf{Skalabilitas.} Pengujian dilakukan dengan jumlah anggota terbatas (5--10 akun pengujian). Analisis skalabilitas untuk jumlah anggota dalam orde ribuan atau puluhan ribu---termasuk strategi \textit{sharding} atau \textit{Layer-2 scaling}---berada di luar lingkup.
\end{enumerate}

\section{Kontribusi Penelitian}

Penelitian ini memberikan empat kontribusi utama:

\begin{enumerate}
	\item \textbf{Implementasi Komprehensif Pertama UU 25/1992 sebagai \textit{Smart Contract}.} Lakomi merupakan sistem pertama yang secara sistematis memetakan 26 ketentuan operasional UU 25/1992 ke dalam empat \textit{smart contract} yang dapat dieksekusi dan diverifikasi. Pendekatan ini berbeda dari penelitian sebelumnya yang bersifat konseptual atau parsial.
	
	\item \textbf{Arsitektur \textit{Dual-Track} untuk Koperasi Blockchain.} Arsitektur Lakomi memisahkan jalur demokrasi (satu-anggota-satu-suara) dari jalur kontribusi finansial (\textit{tier} LTV berbasis simpanan). Pemisahan ini menjawab permasalahan fundamental DAO \textit{token-weighted} yang cenderung plutokratis, sekaligus memberikan insentif finansial bagi anggota untuk meningkatkan kontribusi tanpa mengorbankan kesetaraan suara.
	
	\item \textbf{Metodologi Pemetaan Regulasi-ke-\textit{Smart Contract}.} Penelitian ini menghasilkan metodologi sistematis untuk menerjemahkan ketentuan hukum ke dalam fungsi \textit{smart contract} yang dapat diverifikasi. Metodologi ini bersifat \textit{generalizable} dan dapat direplikasi untuk regulasi lain---misalnya, UU Perseroan Terbatas, UU Yayasan, atau standar ISO.
	
	\item \textbf{Distribusi SHU Enam Kategori On-Chain.} Sistem Lakomi mengimplementasikan distribusi SHU enam kategori sesuai Pasal 45(2) UU 25/1992 secara on-chain untuk pertama kalinya. Parameter distribusi menggunakan \textit{basis points} yang dapat dikonfigurasi melalui tata kelola, memberikan fleksibilitas bagi setiap koperasi tanpa memerlukan \textit{redeployment} kontrak.
\end{enumerate}

\section{Manfaat Penelitian}

Penelitian ini diharapkan memberikan manfaat pada tiga aspek:

\begin{enumerate}
	\item \textbf{Manfaat Akademis.} (a) Memberikan kontribusi pertama berupa implementasi UU 25/1992 sebagai \textit{smart contract} yang lengkap dan terverifikasi; (b) Menghasilkan metodologi pemetaan regulasi-ke-kontrak yang dapat digunakan sebagai acuan untuk penelitian serupa; (c) Memperkaya literatur tentang blockchain untuk koperasi di Indonesia, yang saat ini masih terbatas pada kajian konseptual.
	
	\item \textbf{Manfaat Praktis.} (a) Menyediakan cetak biru (\textit{blueprint}) digitalisasi koperasi yang dapat diadopsi oleh koperasi di Indonesia untuk meningkatkan transparansi dan akuntabilitas; (b) Mendukung fungsi pengawasan koperasi melalui akses \textit{real-time} pengawas terhadap laporan keuangan on-chain; (c) Mengurangi biaya operasional koperasi melalui otomatisasi proses administrasi---pendaftaran anggota, pencatatan simpanan, persetujuan pinjaman, dan distribusi SHU.
	
	\item \textbf{Manfaat bagi Industri dan Pemerintah.} (a) Mendukung program pemerintah ``Koperasi Digital'' dengan menyediakan infrastruktur teknologi yang siap diadopsi; (b) Memanfaatkan infrastruktur DChain sebagai blockchain konsorsium pendidikan tinggi, memperkuat ekosistem blockchain nasional; (c) Menjadi \textit{proof of concept} bahwa regulasi Indonesia dapat diimplementasikan sebagai kode yang dapat dieksekusi, membuka jalan bagi digitalisasi regulasi di sektor lain.
\end{enumerate}

\section{Sistematika Penulisan}

Penulisan skripsi ini disusun dalam enam bab sebagai berikut:

\begin{enumerate}
	\item \textbf{BAB I: Pendahuluan.} Bab ini menguraikan latar belakang penelitian, rumusan masalah, tujuan penelitian, batasan penelitian, kontribusi penelitian, manfaat penelitian, dan sistematika penulisan.
	
	\item \textbf{BAB II: Tinjauan Pustaka dan Dasar Teori.} Bab ini membahas landasan teori yang mendasari penelitian, meliputi: teknologi blockchain dan \textit{smart contract}, \textit{Decentralized Autonomous Organization} (DAO), koperasi di Indonesia, regulasi perkoperasian, dan penelitian terdahulu yang relevan.
	
	\item \textbf{BAB III: Metode Penelitian.} Bab ini menjelaskan metodologi \textit{Design Science Research} yang digunakan, alat dan bahan, alur penelitian, arsitektur sistem, desain \textit{smart contract}, desain \textit{frontend}, alur kerja sistem, dan strategi pengujian.
	
	\item \textbf{BAB IV: Hasil dan Pembahasan.} Bab ini menyajikan hasil implementasi \textit{smart contract} dan \textit{frontend}, hasil pengujian kepatuhan terhadap 26 ketentuan UU 25/1992, analisis biaya gas, perbandingan sistem, dan validasi formula.
	
	\item \textbf{BAB V: Kesimpulan dan Saran.} Bab ini menyajikan kesimpulan yang menjawab rumusan masalah dan saran untuk penelitian selanjutnya.
\end{enumerate}
